\capitulo*{Resumen}
\addcontentsline{toc}{capitulo}{Resumen} 

El presente Trabajo de Fin de Grado, titulado \textit{«Herramienta de gestión para el seguimiento de pacientes con vídeo a través de dispositivo móvil»} (proyecto GII 24.36), se desarrolla en el marco de una colaboración con el departamento de neurología del Hospital Universitario de Burgos. El objetivo principal es el desarrollo de una aplicación móvil nativa para Android, diseñada para facilitar la telerrehabilitación asíncrona entre terapeutas y pacientes. La solución propuesta implementa una arquitectura cliente-servidor que permite a los profesionales sanitarios diseñar rutinas de ejercicios personalizadas y evaluar el desempeño de los pacientes de forma remota.

Un aspecto crítico del proyecto ha sido el diseño de la experiencia de usuario, adaptando la interfaz a las limitaciones motoras de los pacientes mediante el uso de comandos simplificados y flujos guiados. Asimismo, se ha implementado un módulo de captura de vídeo, resolviendo desafíos técnicos complejos relacionados con el ciclo de vida de la cámara y la gestión de archivos locales. Esta captura de vídeo de alta calidad es fundamental, ya que constituirá la base de datos necesaria para futuros módulos de Inteligencia Artificial destinados a la detección de poses y la corrección automática de ejercicios.

El resultado final es una herramienta funcional que permite gestionar roles de usuario (administrador, terapeuta, paciente), realizar y grabar sesiones de rehabilitación, y visualizar la evolución del paciente mediante gráficas de rendimiento, sentando las bases tecnológicas para un sistema de telerrehabilitación inteligente y autónomo.

\capitulo*{Abstract}
\addcontentsline{toc}{capitulo}{Abstract} 

This Bachelor's Thesis, titled \textit{"Management tool for patient follow-up with video via mobile device"} (project GII 24.36), is developed within the framework of a collaboration with the Neurology Department of the University Hospital of Burgos. The main objective is the development of a native Android mobile application designed to facilitate asynchronous telerehabilitation between therapists and patients. The proposed solution implements a client-server architecture that allows healthcare professionals to design personalized exercise routines and assess patient performance remotely.

A critical aspect of the project has been the User Experience (UX/UI) design, adapting the interface to the motor limitations of patients through the use of simplified commands and guided flows. Furthermore, a video capture module has been implemented, solving complex technical challenges related to the camera lifecycle and local file management. This high-quality video capture is fundamental, as it will constitute the necessary database for future Artificial Intelligence modules aimed at pose detection and automatic exercise correction.

The final result is a functional tool that allows for the management of user roles (administrator, therapist, patient), the performance and recording of rehabilitation sessions, and the visualization of patient progress through performance charts, laying the technological foundations for an intelligent and autonomous telerehabilitation system.